\section{Reduced total-gradient implementation for an excited state}
\label{sec:lean}
The separate stationary functionals are useful for derivation and diagnostic
component checks. They do not require two complete production response and
density calculations when only the total gradient is requested. We now
reduce their sum algebraically, keeping the same full moment energy,
independent-energy denominator response, and reference+$P$ eigenvectors.
This section describes a verified algebraic route; it has not replaced the
current oracle and is not a measured production speedup.

\subsection{Projected correction responses}
For a fixed isolated excited root, suppress $\mu$ on the vectors and set
\begin{equation}
 \mathsf B_\mu=\mathsf A-E_\mu\mathbf1,\qquad
 g_r=\bH_{Q,\{0,P\}}^{\T}\bm v,\qquad
 g_l=\bH_{\{0,P\},Q}\bm u.
 \label{eq:leanB}
\end{equation}
The operator $\mathsf B_\mu$ is the shifted projected EOM operator; it is
not the ground CC residual Jacobian J, and is not the ground-state B in
Eq.~\eqref{eq:JBe}. All state and response vectors in this section have
$1+|P|$ components, including the reference component.

The correction is separately linear in $\bm l$ and $\bm r$ at fixed
$T,E,f$, hence
\begin{equation}
 \bm l^{\T}g_l=\bm r^{\T}g_r=\delta_\mu.
 \label{eq:leansolvability}
\end{equation}
Define right and left projected responses by
\begin{align}
 \mathsf B_\mu\bm x&=-(g_l-\delta_\mu\bm r),
       &\bm l^{\T}\bm x&=0,\\
 \mathsf B_\mu^{\T}\bm y&=-(g_r-\delta_\mu\bm l),
       &\bm r^{\T}\bm y&=0.
 \label{eq:leanresponses}
\end{align}
Equation~\eqref{eq:leansolvability} guarantees that the first right-hand
side is orthogonal to the left null vector and the second to the right
null vector. An ordinary inverse of $\mathsf B_\mu$ is inappropriate:
its null vector is the target eigenstate. Use a projected or bordered solve.
For example, define
\begin{equation}
 \mathsf K_\mu=\begin{pmatrix}\mathsf B_\mu&\bm r\\
                                  \bm l^{\T}&0\end{pmatrix}.
\end{equation}
Then the two responses follow from the same bordered operator and its
transpose:
\begin{align}
 \mathsf K_\mu\begin{pmatrix}\bm x\\\xi\end{pmatrix}
 &=\begin{pmatrix}-(g_l-\delta_\mu\bm r)\\0\end{pmatrix},\\
 \mathsf K_\mu^{\T}\begin{pmatrix}\bm y\\\eta\end{pmatrix}
 &=\begin{pmatrix}-(g_r-\delta_\mu\bm l)\\0\end{pmatrix}.
 \label{eq:leanborders}
\end{align}
At a consistent converged solution $\xi=\eta=0$. The borders impose the
normalization-side conditions without requiring a dense projected matrix.
In a selected-tensor implementation, operator actions and preconditioning
can reuse the corresponding EOM right/left machinery.

\subsection{Recovery of the original adjoints and gauge}
For the total objective $F=F_b+F_c$, the on-shell state sources are
\begin{equation}
 F_r=E_\mu\bm l+g_r,\quad F_l=E_\mu\bm r+g_l,
 \quad F_E=-W_Q,
 \quad \nu_{\rm tot}=-E_\mu-\delta_\mu.
 \label{eq:leantotalsources}
\end{equation}
A convenient representative of the original adjoints is
\begin{equation}
 \bm a=\bm y-W_Q\bm l,\qquad \bm b=\bm x.
 \label{eq:leanadjoints}
\end{equation}
Substitution into Eqs.~\eqref{eq:adjr}--\eqref{eq:adjenergy} gives
\begin{align}
 \mathsf B_\mu^{\T}\bm a+F_r+\nu_{\rm tot}\bm l
 &=\mathsf B_\mu^{\T}\bm y+g_r-\delta_\mu\bm l=0,\\
 \mathsf B_\mu\bm b+F_l+\nu_{\rm tot}\bm r
 &=\mathsf B_\mu\bm x+g_l-\delta_\mu\bm r=0,\\
 \bm a^{\T}\bm r+\bm b^{\T}\bm l&=-W_Q.
 \label{eq:leanadjointproof}
\end{align}
The last equality uses $\bm r^{\T}\bm y=\bm l^{\T}\bm x=0$ and
$\bm l^{\T}\bm r=1$. It preserves the denominator's independent-energy
response exactly; that response has not been frozen or discarded.

The representative in Eq.~\eqref{eq:leanadjoints} need not have $b_q=0$.
To recover the original pivot convention, set
\begin{equation}
 \theta=\frac{x_q}{r_q},\qquad
 \bm a'=\bm a+\theta\bm l,\qquad
 \bm b'=\bm b-\theta\bm r.
 \label{eq:leanpivot}
\end{equation}
Then $b'_q=0$, the energy equation is unchanged, and
\begin{equation}
 \bm a'\bm r^{\T}+\bm l(\bm b')^{\T}
 =\bm a\bm r^{\T}+\bm l\bm b^{\T}.
\end{equation}
The corresponding cluster source is invariant for the same reason. Thus
the projected-response gauge and the original constrained pivot gauge give
the identical stationary Hbar coefficient and derivative.

For the base objective alone, Eq.~\eqref{eq:baseeomsimplified} already
provides $\bm a_b=\bm b_b=0$, $\nu_b=-E_\mu$. There is no need to
numerically solve a generic base-state bordered response problem.

\subsection{One combined cluster adjoint and three density factors}
Embed all state/response vectors on reference+$P$ and the divided moments
on Q. Add the base and correction coefficients in
Eqs.~\eqref{eq:excrawcoeff} and \eqref{eq:exccoeff}, then use
Eq.~\eqref{eq:leanadjoints}. Before cluster response, the coefficient is
\begin{align}
 \mathcal A_{\rm pre}
 &=\bm l\bm r^{\T}+\bm v_Q\bm r^{\T}+\bm l\bm u_Q^{\T}
   +(\bm y-W_Q\bm l)\bm r^{\T}+\bm l\bm x^{\T}
   +\mathcal D_\mu\\
 &=\big[(1-W_Q)\bm l+\bm y+\bm v_Q\big]\bm r^{\T}
                 +\bm l(\bm x+\bm u_Q)^{\T}+\mathcal D_\mu.
 \label{eq:leanpre}
\end{align}
Here
\begin{equation}
 \mathcal D_\mu^{\rm full}=\sum_{K\in Q}w_K\ket K\bra K,
 \qquad \mathcal D_\mu^{\rm mp}=W_Q\ket0\bra0.
 \label{eq:leandiagonal}
\end{equation}
The factor $1-W_Q$ is the chain-rule contribution from the uncorrected
state energy in the denominator. It is not an additional energy
renormalization and does not change Eq.~\eqref{eq:energy}.

A single combined cluster multiplier now solves
\begin{equation}
 \boxed{J^{\T}\bm z_t=-\bm q_t,\qquad
 (q_t)_I=\sum_{AB}(\mathcal A_{\rm pre})_{AB}H^I_{AB}.}
 \label{eq:leancluster}
\end{equation}
All state vectors, divided moments, weights and adjoints in this contraction
are held fixed. Their implicit response has already been eliminated by
stationarity. The final coefficient is
\begin{equation}
 \boxed{\mathcal A_{\rm tot}
 =\big[(1-W_Q)\bm l+\bm y+\bm v_Q\big]\bm r^{\T}
 +\bm l(\bm x+\bm u_Q)^{\T}+\bm z_t\evec^{\T}+\mathcal D_\mu.}
 \label{eq:leanthreefactors}
\end{equation}
There are three outer products plus the separate denominator diagonal.
Their pullback and transition-density contractions are those of
Sec.~\ref{sec:densities}. Multiplying two effective ``corrected'' state
vectors together would introduce cross terms absent from
Eq.~\eqref{eq:leanthreefactors}; that is not this reduction.
For MP, the explicit $d_p=-\sum_Kw_K\chi_{Kp}$ Fock density in
Sec.~\ref{sec:fock} must still be added.

\subsection{Production work sequence and reuse}
For a total gradient, the resulting sequence at a fixed geometry is:
\begin{enumerate}
 \item Reuse converged T, Hbar and the target right/left eigenstate from the
       uncorrected energy calculation. Preserve the reference component.
 \item Construct Q moments, denominators and response sources in bulk.
       No response solve is required per Q determinant or Cartesian coordinate.
 \item Solve the two projected state equations in Eq.~\eqref{eq:leanresponses}.
 \item Form the merged source and solve the single cluster equation in
       Eq.~\eqref{eq:leancluster}.
 \item Accumulate one total density from the three factors and denominator
       terms, then solve one total orbital response and contract one set of
       AO derivatives, including overlap and fitted-metric terms.
\end{enumerate}
For fixed SD/selected-T ranks, a production realization should evaluate these
objects through connected tensor or selected-list kernels. Full-CI support
and dense full-space coefficient matrices are not required by the equations.
An exact asymptotic tensor scaling or wall-time speedup requires a separate
contraction/implementation audit.

At a fixed geometry and selection, T, Hbar, J and the HF response operator
can be shared across roots and denominator variants. R/L are shared across
denominator variants for a given root. Different denominators change the
correction sources, so their response solutions are not interchangeable.
The shifted/projected EOM operator also remains root-specific, although its
underlying action and preconditioning structure can be reused.

If separate gradient components are needed, compute the base response once
and cache it. Its state adjoint is analytic; only its cluster response and
density need to be evaluated. Otherwise, the combined route avoids separate
base/correction cluster solves, density builds and orbital-response solves.
This changes the implementation schedule, not the two-functional theory.

\subsection{Algebraic verification and current implementation status}
\label{sec:leanvalidation}
The prototype compares the merged total bare-Hamiltonian gradient
coefficient with the independent
dense constrained oracle for closed/nonclosed P, two excited roots, and
full/MP denominators: eight cases. The maximum coefficient discrepancy is
$\ee{5.551}{-15}$. The checker and saved results are:
\begin{quote}\small
\path{devtools/check_ccpq_lean_response.py}\\
\path{docs/theory/ccpq_lean_response_checks.json}.
\end{quote}
The explicit MP Fock density is retained separately; coefficient agreement
at fixed Fock must not be misreported as permission to omit it.

The original molecular tables continue to document the existing oracle's
validated gradients. They are not timings or independent molecular tests
of a newly optimized implementation. The reduction above is algebraically
verified, but has not replaced the production/oracle path; no measured
speedup or globally minimal implementation is claimed. A production change
should retain the original route as a correctness reference, compare the
resulting densities and nuclear gradients, and record phase-resolved timings.
